%% ============================================================================
%% Elsevier CAS single-column manuscript
%% Transport-inspired population control for rare-event sampling in spin systems
%%
%% Build: pdflatex -> bibtex (not needed; thebibliography is inline) -> pdflatex x2
%% Requires: cas-sc.cls (Elsevier CAS bundle), pgfplots >= 1.16
%% ============================================================================

\documentclass[a4paper,fleqn]{cas-sc}

\usepackage[numbers,sort&compress]{natbib}
\usepackage{amsmath}
\usepackage{amssymb}
\usepackage{booktabs}
\usepackage{graphicx}
\usepackage{caption}
\usepackage{float}
\usepackage{placeins}
\usepackage{enumitem}
\usepackage{pgfplots}
\pgfplotsset{compat=1.16}

\setlist{itemsep=2pt,topsep=3pt}
\newcommand{\WE}{\mathrm{WE}}
\newcommand{\FOM}{\mathrm{FOM}}

%% ---------------------------------------------------------------------------
%% FIGURE 1 DATA.
%% The rows below are ONLY the values quoted in the text. Replace this block
%% with the full binned tail from the production output
%% (columns: |m|  P_naive  P_we ; use "nan" where a method records no samples).
%% ---------------------------------------------------------------------------


\begin{document}
%% cas-sc redefines the figure environment after float is loaded; restore [H]
\restylefloat{figure}
\restylefloat{table}
\let\WriteBookmarks\relax
\def\floatpagepagefraction{1}
\def\textpagefraction{.001}

\shorttitle{Transport-inspired rare-event sampling in spin systems}
\shortauthors{M.Y. Alzaatreh}

\title[mode=title]{Transport-Inspired Rare-Event Sampling in Spin Systems: Learned
Importance Coordinates Improve Access to Deep Targets}

\author{M.Y. Alzaatreh}
\cormark[1]
\ead{mzaatreh@fbsu.edu.sa}
\credit{Conceptualization, Methodology, Software, Formal analysis, Investigation,
Visualization, Writing -- original draft, Writing -- review \& editing}

\affiliation{organization={Natural Sciences Unit, Fahad Bin Sultan University},
            city={Tabuk},
            country={Saudi Arabia}}

\cortext[1]{Corresponding author}

% ---- Abstract --------------------------------------------------------------
\begin{abstract}
This work transfers the variance-reduction toolkit of radiation transport, adjoint importance
functions and weight-window population control, to rare-event sampling in lattice spin systems,
using the shared backward-conditioning structure of the transport adjoint and the transition-path
committor. Exact references validate the implementation: a birth--death process reproduces the
zero-variance limit of exact importance, and enumeration of a $4\times4$ lattice confirms that
weight-conserving Weighted Ensemble (WE) resampling is unbiased for any binning coordinate. At
matched spin-update cost, WE reaches an $L=20$ Ising magnetization-tail bin about $24$ times below
the direct-sampling floor of $5.9\times10^{-6}$. For the frustrated Edwards--Anderson spin glass,
the weighted population is binned by a self-consistency-trained neural coordinate acting on the
full spin configuration. At the deepest calibrated target at $L=16$, pooled across three
independent disorder realizations with one trained network each, the event is observed in
$63/90$ learned-coordinate replicas against $46/90$ for the physical energy coordinate (Fisher
exact $p=0.014$); a complementary check pooling three independently trained networks on a single
realization against a matched-size baseline gives $60/90$ against $44/90$ ($p=0.023$). The gain
is in reliability of access, not efficiency: repeated network inference costs four to five times
the wall-clock time of the energy coordinate, leaving a directly measured, two-level-bootstrapped
figure of merit at roughly one-third of parity ($0.33$, $95\%$ CI $[0.19,0.58]$). The transport
toolkit therefore transfers, the reliability advantage generalizes across independent disorder
realizations rather than being tied to one lattice, and the remaining barrier is inference cost
rather than coordinate quality.
\end{abstract}

% ---- Highlights ------------------------------------------------------------
\begin{highlights}
\item Adjoint importance and the rare-event committor share the same backward-conditioning
structure for first-passage problems, and coincide for an absorbing target.
\item Exact tests validate both the zero-variance importance limit and unbiased WE
population control.
\item At matched update cost, WE reaches a tail bin approximately $24$ times below the naive
sampling floor near $5.9\times10^{-6}$, where direct sampling returns no estimate at all.
\item A learned importance coordinate significantly improves deep-event discovery in the
Edwards--Anderson spin glass, pooled across three independent disorder realizations
($70.0\%$ versus $51.1\%$, $p=0.014$) and, separately, across three trained networks on one
realization ($66.7\%$ versus $48.9\%$, $p=0.023$).
\item The improvement is in reliability of access, not efficiency: directly measured inference
cost leaves the learned coordinate at roughly one-third of parity with the energy coordinate in
wall-clock figure of merit ($0.33$, $95\%$ CI $[0.19,0.58]$).
\end{highlights}

% ---- Keywords --------------------------------------------------------------
\begin{keywords}
Rare-event sampling \sep committor \sep adjoint importance \sep Weighted Ensemble \sep
variance reduction \sep Ising model \sep spin glass \sep neural importance
\end{keywords}

\maketitle

% ============================================================================
\section{Introduction}\label{sec:intro}

A rare event is one that direct simulation almost never observes: the nucleation of a new
phase, a protein crossing between conformational basins, or a detector registering a
particle that has penetrated many mean free paths of shielding. In each case the probability
of the event is itself the quantity of scientific interest, and it is exactly that quantity
which unbiased dynamics represents worst. If an event has equilibrium or first-passage
probability $p\ll1$, the relative error of a direct Bernoulli estimator formed from $N$
effectively independent samples follows the elementary result \citep{kahn1951,kalos2008}
\begin{equation}
\varepsilon_{\rm rel}=\sqrt{\frac{1-p}{Np}}\;\simeq\;(Np)^{-1/2},\qquad p\ll1,
\label{eq:bernoulli}
\end{equation}
so the sample count required for fixed relative precision grows as $1/p$ and each additional
decade of rarity costs an order of magnitude in computation. The failure is
therefore structural rather than incidental, and every established remedy responds to it in
the same way: not with more samples, but with better-placed ones. This elementary scaling is
the origin of a broad family of rare-event techniques, including importance sampling,
splitting, transition-path sampling, Forward Flux Sampling (FFS), Adaptive Multilevel
Splitting (AMS), and Weighted Ensemble (WE) methods \citep{kahn1951,kalos2008,huber1996,
zuckerman2017,allen2009,cerou2007,bolhuis2002}. What each of them requires in practice is a
rule for deciding which samples are the better-placed ones.

Two mature traditions answer that question using closely related mathematical objects. In
transition-path theory, the committor $h(s)$ is the probability that a trajectory initiated
at configuration $s$ reaches a target set before a competing boundary, and it provides an
optimal reaction coordinate for that event \citep{e2006,metzner2009,vandeneijnden2010}. In
classical particle transport, an adjoint solution assigns an importance to a phase-space
point according to its expected contribution to a detector response; this importance can be
used to bias sources, alter transition probabilities, and construct weight windows
\citep{lewins1965,spanier1969,bell1970,wagner1998,haghighat2003}. The formal connection to
the Doob $h$-transform and conditioned Markov processes is well known
\citep{doob1984,touchette2009,chetrite2015}. The contribution of the present work is
therefore not to claim a new identity, but to transfer the associated \emph{transport
variance-reduction toolkit} to lattice spin dynamics and test, under controlled validation,
which parts remain effective.

The transport perspective is attractive because it separates the ideal importance function
from the population-control mechanism used when that importance is only approximate. Exact
adjoint biasing can in principle produce zero-variance estimators, while weight windows use
splitting and roulette to keep statistical weight near a prescribed range. This distinction
underlies CADIS and related hybrid deterministic--Monte Carlo variance-reduction methods
\citep{wagner1998,haghighat2003,wagnerfw2014}. In the present setting, the same logic
motivates using a committor-like quantity to organize a weighted population while retaining
the original spin dynamics.

Rare-event population methods have a parallel history outside transport. WE was introduced
for biomolecular transition sampling by Huber and Kim, shown to be statistically exact for
broad classes of dynamics and binning procedures, and made practical for complex molecular
systems by steady-state formulations and mature software
\citep{huber1996,zhang2010,bhatt2010,russo2022}; FFS and AMS address related first-passage
problems through interfaces or adaptive levels \citep{allen2009,valeriani2007,cerou2007},
while transition-path sampling works with reactive-path ensembles and can avoid prescribing a
one-dimensional coordinate at all \citep{dellago1998,bolhuis2002}. Machine learning
meanwhile offers several routes to the high-dimensional objects these methods need: flow-based
and autoregressive samplers for equilibrium distributions
\citep{wu2019,noe2019,nicoli2020,gabrie2022}, neural approximations of committor and
Feynman--Kac observables from short trajectories \citep{khoo2019,strahan2023}, learned bias
potentials and importance functions \citep{hua2024,kimcai2026}, and neural Doob dynamics for
large deviations \citep{zhao2025}. Together these motivate the central practical question of
this paper: when the rare event occurs in a high-dimensional spin configuration space, how much
is gained by replacing a hand-designed coordinate with a learned, committor-inspired one --
measured against a strong physical baseline and with inference cost charged to the result?

The study is therefore organized as a validation ladder rather than a single benchmark. A
solvable one-dimensional birth--death chain first verifies the zero-variance limit of exact
importance; exact enumeration of a $4\times4$ Ising lattice then verifies unbiased WE
resampling; a larger $L=20$ Ising calculation tests whether WE genuinely extends a rare
magnetization tail at matched spin-update cost; and finally the Edwards--Anderson (EA) spin
glass provides a frustrated landscape in which magnetization is a poor coordinate and the
low-energy target is a more demanding test for learned configuration-space importance
\citep{ising1925,onsager1944,edwards1975,binder1986}. The hierarchy is designed to separate
three claims: the exact mathematical importance limit, unbiased weighted population control,
and the ability of a learned coordinate to improve access to a genuinely rare target.
Throughout, \emph{reliability} (whether a method reaches the event at all, reproducibly) is
distinguished from \emph{efficiency} (variance reduction per unit of actual computational
cost); the distinction is essential for neural coordinates, because inference overhead is not
represented by a spin-update count.

The main findings are positive but bounded. WE is unbiased and extends the accessible Ising
tail beyond the direct-sampling floor at matched update cost. In the EA model, a
self-consistency-trained neural coordinate significantly improves deep-event discovery over
the physical energy coordinate, an advantage that is reproduced across independent disorder
realizations rather than tied to one lattice. That improvement is a reliability result: once
inference is included in the cost, the same coordinate does not improve the figure of merit.
Together,
these results identify a transferable theoretical core, a validated route for using learned
importance information inside a weight-conserving rare-event sampler, and the specific
barrier that a practical implementation must still clear.

\FloatBarrier

% ============================================================================
\section{Theory and Methods}\label{sec:methods}

\subsection{Adjoint importance, committors, and zero-variance sampling}\label{sec:adjoint}
Let $K(s,s')$ be the transition kernel of a Markov chain and let $B$ denote a target set.
For a first-passage event, the forward committor is the probability of reaching $B$ before
the competing set $A$ \citep{e2006,metzner2009}:
\begin{equation}
h(s)=\Pr_{s}\!\left(\tau_B<\tau_A\right).
\label{eq:committor}
\end{equation}
For a discrete Markov process, this committor satisfies the backward (harmonic) equation in
the interior of the state space, a standard result of transition-path theory for Markov jump
processes \citep{e2006,metzner2009,vandeneijnden2010}:
\begin{equation}
h(s)=\sum_{s'}K(s,s')\,h(s'),\qquad s\notin A\cup B,
\label{eq:harmonic}
\end{equation}
with boundary values $h=0$ on $A$ and $h=1$ on $B$. Conditioning the dynamics with a
positive harmonic function leads to the Doob $h$-transform, introduced in classical potential
theory \citep{doob1984} and used in this form throughout the large-deviation and
conditioned-process literature \citep{touchette2009,chetrite2015}:
\begin{equation}
K_h(s,s')=K(s,s')\,\frac{h(s')}{h(s)},
\label{eq:doob}
\end{equation}
where defined. This is the Markov-chain counterpart of adjoint-based importance biasing in
transport. With the exact $h$, successful paths acquire identical total statistical weight,
producing the familiar zero-variance ideal.

The relationship between the two importance functions should be stated precisely, since it is
weaker than an identity. The committor of Eq.~\eqref{eq:committor} is a probability: the chance
of reaching $B$ before $A$ from $s$, bounded in $[0,1]$ and defined only once a competing
boundary is specified. The transport adjoint is an expected contribution to a specified
detector response, carries the units of that response, and is unbounded in general. What the
two share is structure rather than definition: both are solutions of a backward equation
generated by the same transition operator, both express the value of a state with respect to a
terminal objective, and both therefore play the same role when used to tilt or to bin. They
coincide as functions in the particular case of an absorbing target with unit detector
cross-section and no multiplication, which is the case constructed here; they are not
universally the same object, and this paper does not claim they are. What transfers is the
backward-conditioning structure and the machinery built on it, not an equivalence of the two
quantities in general.

This identity motivates two distinct uses of approximate importance. The first changes
transition probabilities directly, which can be fragile when $h$ is inaccurate. The second
retains the original dynamics but uses population splitting and merging to control
statistical weight. The second strategy is adopted here because weight conservation
preserves unbiasedness even when the binning coordinate is imperfect.

\subsection{Zero-variance path weighting and the transport correspondence}\label{sec:zerovar}
The zero-variance statement can be seen directly at the path level. Consider a trajectory
$\omega=(s_0,s_1,\dots,s_\tau)$ generated under the Doob-transformed kernel until it reaches
either $A$ or $B$. Substituting the Doob transform of Eq.~\eqref{eq:doob} into the path
likelihood ratio gives the telescoping identity below, derived here for the discrete
first-passage setting from Eqs.~\eqref{eq:harmonic}--\eqref{eq:doob}; it is the discrete-state
counterpart of the classical zero-variance principle of transport importance sampling
\citep{lewins1965,spanier1969,haghighat2003} and of the conditioned-process construction of
Doob \citep{doob1984} and Chetrite and Touchette \citep{chetrite2015}:
\begin{equation}
\prod_{t=0}^{\tau-1}\frac{K(s_t,s_{t+1})}{K_h(s_t,s_{t+1})}
=\prod_{t=0}^{\tau-1}\frac{h(s_t)}{h(s_{t+1})}
=\frac{h(s_0)}{h(s_\tau)} .
\label{eq:telescope}
\end{equation}
On a successful path ending in $B$, $h(s_\tau)=1$, so every successful history contributes
exactly $h(s_0)$.

The resulting estimator should be stated explicitly, since the zero-variance property depends
on how it is constructed. Fix an initial state $s_0$ with $h(s_0)>0$ and let the target
quantity be the first-passage probability $\pi=h(s_0)=\Pr_{s_0}(\tau_B<\tau_A)$. Generate $M$
independent trajectories under $K_h$ and assign each the likelihood-ratio weight of
Eq.~\eqref{eq:telescope}, giving
\begin{equation}
\hat\pi=\frac1M\sum_{m=1}^{M}
\left[\prod_{t=0}^{\tau_m-1}\frac{K(s_t^{(m)},s_{t+1}^{(m)})}{K_h(s_t^{(m)},s_{t+1}^{(m)})}\right]
\mathbf 1_B\!\left(s_{\tau_m}^{(m)}\right).
\label{eq:zvest}
\end{equation}
Two properties of $K_h$ make this exact. States with $h=0$ receive zero transition probability,
so the transformed chain never enters $A$ and every trajectory terminates in $B$; and on each
such trajectory the weight telescopes to $h(s_0)$ by Eq.~\eqref{eq:telescope}. Every term of
the sum therefore equals $h(s_0)$, so $\hat\pi=\pi$ identically and
$\mathrm{Var}(\hat\pi)=0$ for any $M\ge1$. The statement is specific to a first-passage
probability from a fixed initial state, requires $h(s_0)>0$, and requires the exact $h$; it
degrades continuously, and in general non-uniformly, when $h$ is replaced by an approximation.
This is the property verified numerically in \S\ref{sec:bd}.

For an approximate importance function $\tilde h$, the exact telescoping property is no
longer available. The present implementation therefore separates importance representation
from dynamical biasing: the physical Metropolis kernel is retained, while the learned or
physical progress coordinate is used only to organize population splitting and merging. This
design allows approximate importance information to guide computational effort while weight
conservation preserves the target measure.

Table~\ref{tab:map} summarizes the operational correspondence used throughout the paper. The
mapping is not intended to imply that all transport constructs carry over unchanged. Source
biasing, for example, has no literal analogue when the equilibrium spin distribution is fixed
and the goal is to sample a stationary tail. The useful transfer is instead structural: an
adjoint-like scalar identifies configurations that matter to the response, and a
weight-management mechanism concentrates trajectory population in those configurations.

\begin{table}[H]
\centering\small
\caption{Transport concepts and their rare-event spin-system counterparts used in this work.
The final column states the role tested here.}
\label{tab:map}
\begin{tabular}{p{0.26\textwidth}p{0.28\textwidth}p{0.34\textwidth}}
\toprule
Radiation-transport concept & Markov/spin counterpart & Role in this study \\
\midrule
Adjoint importance & Committor / event-conditioned value function & Defines ideal progress
toward a rare target \\
CADIS importance bias & Doob $h$-transform / tilted dynamics & Verified in the solvable
zero-variance demonstration \\
Weight windows & Weighted splitting and merging & Implemented through WE resampling with
conserved statistical weight \\
Detector response & Rare-event indicator or threshold probability & Magnetization-tail or
low-energy probability \\
\bottomrule
\end{tabular}
\end{table}

\subsection{Spin models and rare-event targets}\label{sec:models}
For the nearest-neighbor ferromagnetic Ising model, the Hamiltonian is that introduced by
Ising \citep{ising1925} and solved exactly in two dimensions by Onsager \citep{onsager1944},
\begin{equation}
H(\mathbf s)=-J\sum_{\langle ij\rangle}s_is_j,\qquad s_i\in\{-1,+1\},
\label{eq:ising}
\end{equation}
with periodic boundaries and $J=1$. The principal rare observable is the absolute
magnetization $|m|$. For the spin glass, the corresponding random-bond Hamiltonian is that of Edwards and Anderson
\citep{edwards1975}, in the $\pm J$ form reviewed by Binder and Young \citep{binder1986},
\begin{equation}
H(\mathbf s)=-\sum_{\langle ij\rangle}J_{ij}s_is_j,\qquad J_{ij}\in\{-1,+1\},
\label{eq:ea}
\end{equation}
with one fixed disorder realization per comparison. Quenched disorder and frustration make
the low-energy states disordered, so $m\approx0$ irrespective of energy and magnetization is
not an adequate progress coordinate; low energy is used as the rare-event target instead. All
reported production comparisons use single-spin Metropolis dynamics \citep{metropolis1953}
at $T=2.6$ unless stated otherwise.

\subsection{Weighted Ensemble estimator and resampling}\label{sec:we}
WE propagates a weighted population under the original, unbiased Metropolis dynamics for a
resampling interval $\tau$. After propagation, walkers are partitioned into bins defined by a
coordinate $q(s)$. If the total statistical weight in bin $b$ is $W_b$, resampling
redistributes that same $W_b$ among the target number of walkers in the bin. Splitting
therefore creates additional trajectories with reduced individual weights, while merging
removes redundant trajectories without changing total weight. Because the dynamics themselves
are not altered and the resampling operation conserves weight, the coordinate controls
efficiency rather than correctness.

For a target set $B$, WE estimates the target probability from the total statistical weight
carried by walkers in $B$. The estimator takes the weight-conserving form established for WE
by Huber and Kim \citep{huber1996}, proved statistically exact for broad classes of dynamics
and binning procedures by Zhang et al.\ \citep{zhang2010}, and used in the steady-state and
review formulations of Bhatt et al.\ \citep{bhatt2010} and Zuckerman and Chong
\citep{zuckerman2017}:
\begin{equation}
\hat\pi_t=\frac{\sum_i w_i^{(t)}\,\mathbf 1_B\!\left(s_i^{(t)}\right)}{\sum_i w_i^{(t)}},
\label{eq:weest}
\end{equation}
and the reported estimate is averaged over post-burn-in resampling cycles. Three coordinates
are compared: absolute magnetization $|m|$, energy per spin $E/N$, and a learned scalar
coordinate $I_\theta(s)$. Magnetization is deliberately retained as a mismatched baseline for
the EA spin glass, whereas energy is a natural physical baseline for the low-energy target.

\subsection{Rarity calibration and experimental protocol}\label{sec:calib}
A fixed rare-event threshold can be either trivial or physically unreachable when lattice
size, temperature, or disorder changes. The computational protocol therefore uses a short
pilot simulation to locate the most extreme value reached under ordinary dynamics and then
places the production threshold beyond that pilot extreme by a small number of discrete
lattice steps. For the Ising magnetization target the lattice spacing is $2/L^2$; for the EA
energy-per-spin target it is $4/L^2$. The integer ``beyond'' used below records how many such
steps the target is displaced past the pilot extreme. This convention provides an operational
depth scale without claiming that equal values of ``beyond'' correspond to identical
probabilities at different system sizes.

The resampling interval is also treated as a physical parameter rather than a bookkeeping
choice. Independent diagnostics in the EA model show rapid relaxation of deep configurations:
after one Metropolis sweep, the mean energy of a set of deep configurations moves
substantially back toward the bulk. Consistent with this, shorter resampling intervals
penetrate the deep ladder more effectively than long intervals at comparable cost. The
production deep-rare calculation therefore uses $\tau=2$, while the exact and shallow
validation calculations retain their own explicitly stated settings.

\subsection{Neural importance coordinate}\label{sec:neural}
The learned coordinate is a compact periodic convolutional network acting on the full spin
configuration. The implementation uses two circularly padded $3\times3$ convolutional layers
with nonlinear activations, followed by global mean and maximum pooling and a scalar output.
This architecture is intentionally small: the research question is whether a
configuration-space coordinate can discover useful progress information that is absent from a
simple scalar observable, not whether a large neural model can fit the training set.

During development, several training targets were tested. Rollout-based dense surrogates were
useful for avoiding the all-zero labels produced by a direct binary committor target in very
rare regimes, but such a label is by construction energy-adjacent, since it records the best
energy reached in a short walk. Later experiments therefore used objectives shaped by the
harmonic self-consistency condition of Eq.~\eqref{eq:harmonic} itself. This direction sits
within a broader neural-committor literature, including direct variational approximation of
high-dimensional committors, semigroup formulations that avoid explicit second derivatives,
active importance sampling for rare-event-dominated objectives, short-trajectory
Feynman--Kac training, and self-consistent committor-enhanced sampling
\citep{khoo2019,li2022,rotskoff2022,strahan2023,kang2024,trizio2025}. The self-consistency
objective used for the discrete MCMC comparison is most directly connected to the recent
neural importance-sampling formulation of Kim and Cai \citep{kimcai2026}, in which a learned
importance function is coupled to biased transitions and branching. Two practical points
follow. First, the self-consistency condition together with the target boundary alone is
under-determined, since the constant solution satisfies it everywhere; a second boundary
condition anchoring bulk-equilibrium states to a distinct reference value is necessary rather
than optional. Second, the present study uses the learned function more conservatively than
Kim and Cai, as a WE binning coordinate, so that coordinate quality can be evaluated without
introducing likelihood-ratio bias from an inaccurate directional tilt.

A terminological point follows from this design and is maintained throughout. The trained
network $I_\theta$ is a \emph{learned importance coordinate} whose objective is shaped by the
harmonic condition of Eq.~\eqref{eq:harmonic}; it is not claimed to be a quantitatively
accurate estimate of the committor $h$, and no such accuracy is measured here. Establishing
that would require a reference committor on a system where one can be computed, which is a
separate exercise from the sampling question addressed in this paper. Because $I_\theta$ enters
only through the binning coordinate, and because weight-conserving resampling is unbiased for
any binning, the results below do not depend on how close $I_\theta$ is to $h$: an inaccurate
coordinate costs efficiency, never correctness. The word ``committor'' is therefore reserved
for the exact object of \S\ref{sec:adjoint}, and $I_\theta$ is described as
committor-inspired rather than as a committor.

\subsection{Computational cost and reporting}\label{sec:cost}
Two cost measures are used, and the distinction between them matters for the interpretation
of the results. Matched-cost comparisons of sampling reach use the number of attempted spin
updates as the common Monte Carlo work unit; this measure is natural for the Ising-tail
comparison because both direct sampling and WE use the same Metropolis propagation kernel.
It is, however, blind to neural inference, which is performed repeatedly during WE binning
and does not appear in a spin-update count. Wall-clock time is therefore also recorded for
the learned-coordinate runs and is the measure used whenever a figure of merit is quoted for
a neural coordinate. The figure of merit is the cost-normalized inverse variance standard in
transport variance-reduction studies \citep{spanier1969,wagnerfw2014},
\begin{equation}
\FOM=\frac{1}{\sigma_{\rm rel}^{2}\,T},
\label{eq:fom}
\end{equation}
in which $T$ is the computational cost charged to the calculation; the two cost measures above
give the two values of $T$ reported in this work.

Independent replicas are used to quantify sampling variability, and zero-hit replicas are
retained explicitly because non-observation is itself meaningful in the deep rare-event
regime. The primary learned-coordinate result is reported through event-discovery frequency
across independent runs, complemented by Fisher's exact test; cost-normalized efficiency is
reported separately in \S\ref{sec:costdisc}.

\subsection{Statistical analysis and uncertainty quantification}\label{sec:stats}
Suppose $R$ independent replicas produce estimates $\hat\pi_r$. Following the independent-replica
uncertainty reporting recommended for WE studies \citep{zuckerman2017} and implemented in
WESTPA \citep{russo2022}, the across-replica relative standard deviation is summarized as
\begin{equation}
\sigma_{\rm rel}=\frac{\mathrm{sd}(\hat\pi_1,\dots,\hat\pi_R)}{\bar\pi},\qquad
\bar\pi=\frac1R\sum_{r=1}^{R}\hat\pi_r .
\label{eq:relsd}
\end{equation}
For deep targets, however, the most informative quantity can be the fraction of independent
runs that observe the event at all. Zero-hit counts are therefore reported explicitly rather
than being hidden inside a variance summary.

For the deep-rare neural comparison, event-positive and zero-hit replicas are compared using
the exact conditional test for a $2\times2$ contingency table introduced by Fisher
\citep{fisher1922}, which is appropriate here because the expected counts are small. Where
confidence intervals on ratios of figures of merit are quoted, they are obtained by the
nonparametric bootstrap \citep{efron1979} in a two-level form that resamples the
trained-network axis in the outer loop and the replica axis of each network in the inner loop.

\FloatBarrier

% ============================================================================
\section{Validation}\label{sec:validation}

\subsection{Exact-importance limit}\label{sec:bd}
A biased one-dimensional birth--death chain provides a solvable test. Starting at $s_0=1$
with absorbing boundaries at $0$ and $N=15$ and step probability $p=0.4$ toward the rare
boundary, the exact committor is $h(s_0)=1.144443\times10^{-3}$. Direct Monte Carlo gives a
relative variance of $2.17\times10^{-3}$ for $4\times10^{5}$ histories, in agreement with the
Bernoulli prediction $1/(N p)$. Using the exact importance transform reduces the measured
relative variance to approximately $10^{-30}$, numerically reproducing the zero-variance
prediction of Eq.~\eqref{eq:telescope}. This exact result provides the reference point for
the weight-conserving population strategy used in the spin-system calculations.

\subsection{Exact-enumeration gate for WE}\label{sec:gate}
For the $4\times4$ Ising model, all $2^{16}$ configurations can be enumerated. Across the
three validation regimes used in the project---above $T_c$, below $T_c$ with the double-well
structure, and the folded $|m|$ distribution---the mean $|\Delta\ln P|$ between WE and exact
enumeration ranges from approximately $0.006$ to $0.024$. The $T=2.6$ case gives the best of
these values, $0.006$, and is retained as the representative detailed gate. A separate
learned-importance gate gives target-probability ratios of $1.000$ and $1.001$ relative to
exact enumeration when WE is binned by magnetization and energy, respectively. These tests
establish the implementation property needed below: resampling is unbiased even when the
coordinate is not optimal.

The purpose of this ladder, summarized in Table~\ref{tab:ladder}, is methodological. The first
test verifies the ideal importance limit, the next two verify unbiased population control
against exact answers, and only then is the method used where an exact distribution is
unavailable. Consequently, later differences between coordinates are interpreted as efficiency
or reliability differences rather than as evidence that one coordinate changes the target
measure.

\begin{table}[H]
\centering\small
\caption{Validation ladder used before interpreting the larger rare-event calculations. Each
row validates a different component of the computational framework.}
\label{tab:ladder}
\begin{tabular}{p{0.30\textwidth}p{0.24\textwidth}p{0.34\textwidth}}
\toprule
Test & Reference & Outcome \\
\midrule
Birth--death importance transform & analytic committor & exact-importance relative variance
$\sim10^{-30}$ \\
$4\times4$ Ising WE distribution & exact enumeration & mean $|\Delta\ln P|=0.006$--$0.024$
across three regimes \\
$4\times4$ learned-coordinate gate & exact target probability & WE ratios $1.000$ and
$1.001$ \\
$L=20$ Ising tail & matched-update naive MC & WE resolves one additional rare-tail bin to
$2.4\times10^{-7}$ \\
\bottomrule
\end{tabular}
\end{table}

\FloatBarrier

% ============================================================================
\section{Rare-Event Reach in the Ising Ferromagnet}\label{sec:tail}

\subsection{WE extends the Ising tail}\label{sec:tailreach}
At $L=20$ and matched update cost ($\sim3.4\times10^{8}$ spin updates), naive sampling and WE
agree throughout the well-sampled region, and their behavior separates only in the far tail.
Naive sampling reaches $|m|=0.886$ at $P\approx5.9\times10^{-6}$ but records no samples in the
next tail bin, whereas WE estimates $P\approx2.4\times10^{-7}$ at $|m|=0.932$. The claim supported by this comparison is that WE reaches a tail bin
approximately $24$ times below the direct-sampling floor under the matched update budget,
although the deepest estimate is itself statistically noisy: it carries an across-replica
relative standard deviation of approximately $2.0$ over four replicas. The robust content is
therefore that WE returns an estimate where direct sampling returns none at all, rather than
that the value $2.4\times10^{-7}$ is precise. Expressed as sampling reach, the lowest probability
resolved under the common budget falls from $5.9\times10^{-6}$ to $2.4\times10^{-7}$, an
extension of approximately $1.4$ orders of magnitude, and this is the quantity the rest of the
paper depends on.



\subsection{Calibration of the deep-rare spin-glass regime}\label{sec:depth}
The neural-coordinate experiment is placed deliberately beyond the range where direct sampling
is consistently informative. The calibrated EA thresholds at $L=16$ become progressively less
accessible under the same direct-sampling budget. Table~\ref{tab:depth} summarizes this depth
calibration and identifies the access-dominated regime used for the learned-coordinate test.

\begin{table}[H]
\centering\small
\caption{Depth calibration for the EA model at $L=16$ using 10 independent direct-sampling
replicas. The integer ``beyond'' indexes successively deeper discrete energy thresholds.}
\label{tab:depth}
\begin{tabular}{lrrr}
\toprule
Beyond & $\hat\pi$ & relative s.d. & zero-hit replicas \\
\midrule
1 & $5.29\times10^{-5}$ & 0.23 & 0/10 \\
2 & $5.48\times10^{-6}$ & 0.46 & 0/10 \\
3 & $7.14\times10^{-7}$ & 1.61 & 7/10 \\
4 & $0$ & --- & 10/10 \\
\bottomrule
\end{tabular}
\end{table}

The beyond$=4$ threshold is therefore used as a stress test of rare-event access. This choice
makes the learned-coordinate comparison a direct test of whether a configuration-space
importance map can organize the weighted population so that trajectories continue to enter a
region that is essentially absent from the corresponding direct Monte Carlo sample.

\FloatBarrier

% ============================================================================
\section{Learned Importance in the Deep-Rare Spin Glass}\label{sec:deep}

At beyond$=4$ ($E/N\le-1.119$), direct Monte Carlo is effectively blind under the calibrated
budget. The WE budget was increased to 80 bins, 40 walkers per occupied bin, 3000 iterations,
and $\tau=2$. Three independently trained self-consistency networks were then compared with
the same physical-coordinate baselines. Using independently trained networks is important
because the objective is to establish a reproducible coordinate-level effect rather than a
favorable outcome from one particular initialization.

\begin{table}[H]
\centering\small
\caption{Deep-rare EA comparison at $L=16$, beyond$=4$. Each learned-coordinate row is a
separately trained network evaluated over 30 replicas. Figures of merit are quoted on the
spin-update cost measure of \S\ref{sec:cost}; the wall-clock correction is applied in
\S\ref{sec:costdisc}.}
\label{tab:deep}
\begin{tabular}{lrrr}
\toprule
Method & event-positive replicas & $\FOM_{\rm MC}$ & ratio vs.\ $\WE[E]$ \\
\midrule
naive & 1/30 & $2.48\times10^{-10}$\textsuperscript{a} & --- \\
$\WE[m]$ & 5/30 & $1.04\times10^{-10}$ & 0.62 \\
$\WE[E]$ & 12/30 & $1.69\times10^{-10}$ & 1.00 \\
$\WE[I_\theta]$, seed 0 & 19/30 & $2.18\times10^{-10}$ & 1.29 \\
$\WE[I_\theta]$, seed 1 & 20/30 & $1.81\times10^{-10}$ & 1.07 \\
$\WE[I_\theta]$, seed 2 & 21/30 & $1.77\times10^{-10}$ & 1.04 \\
\bottomrule
\end{tabular}
\\[2pt]
{\footnotesize\textsuperscript{a}Not a meaningful figure of merit: driven by a single
event-positive replica out of 30 (relative s.d.\ $=5.5$); reported for completeness only.
The underlying $\hat\pi$ estimates for the weight-window methods remain individually noisy in
this regime, with relative s.d.\ between $1.2$ and $1.9$ at 30 replicas.}
\end{table}

The primary evidence is the network-level replication. Each of the three independently
trained networks exceeds the energy coordinate on its own: discovery rates of $63\%$, $67\%$,
and $70\%$ (19/30, 20/30, 21/30) against $40\%$ (12/30) for $\WE[E]$
(Figure~\ref{fig:deeprare}). The effect is therefore
reproduced three times from independent initializations, with no overlap in direction, which is
the property that matters for a coordinate-level claim.

Pooling those runs gives 60 event-positive replicas out of 90 ($66.7\%$) against a naive
reading of 12/30 ($40.0\%$), and Fisher's exact test on the resulting table gives $p=0.017$.
That $\WE[E]$ figure is a single fixed 30-replica draw, however, reused as the comparison
point for all three networks; a further 60 independent replicas (disjoint random seeds, same
configuration) give a better-calibrated estimate of its true discovery rate, $44/90$
($48.9\%$) pooled, close to but above the original draw, which happened to land on the low
side by ordinary sampling variation. Repeating the test against this balanced baseline gives
$66.7\%$ versus $48.9\%$, $p=0.023$ -- a smaller margin than the naive comparison but still
significant, and the value that should be quoted going forward. This pooled test is
offered as supporting evidence rather than as the sole statistical basis, because the 90
learned-coordinate replicas are grouped within three trained networks and are not fully
exchangeable. The grouping appears benign here, since the three per-network rates lie within
$7$ percentage points of one another and a network-level effect large enough to invalidate
pooling would be expected to produce visible heterogeneity; but the replication across networks,
not the pooled $p$-value, is what the conclusion rests on. This is the central within-disorder
learned-coordinate result: the learned importance map increases the probability that a
finite-budget weighted-ensemble run reaches the deep target. \S\ref{sec:disordergen} extends
this replication axis from network initialization to independent disorder realizations.

The comparison is especially informative because energy is a strong physical baseline for this
task. The target itself is defined by low energy, so an improvement over $\WE[E]$ cannot be
attributed to comparing against an irrelevant coordinate. The network instead receives the full
spin configuration and can distinguish configurations with similar instantaneous energy but
different local structure and different dynamical prospects of progressing toward the target.
This is also the mechanistic reason to prefer a self-consistency objective over the rollout
surrogate of \S\ref{sec:neural}: two configurations at identical energy can have very different
committor values on a frustrated landscape, and an energy-adjacent label cannot represent that
difference. The present experiment does not attempt to assign a unique microscopic
interpretation to the learned representation; it establishes that the representation contains
useful progress information beyond the scalar energy coordinate.

The role of the learned function is deliberately conservative. It is used to define the WE
partition rather than to replace the physical transition kernel. Because the split/merge
operation conserves total statistical weight, the estimator remains unbiased for any binning
coordinate. The learned model therefore influences how computational effort is distributed
across configuration space without changing the equilibrium target measure. This separation is
attractive for extending neural rare-event methods to more complex spin systems, where direct
use of an approximate committor inside a likelihood-ratio tilt would require stronger control
of approximation error.

\subsection{Generalization across disorder realizations}\label{sec:disordergen}
The result above establishes reproducibility across network initializations on one $J_{ij}$
sample; it does not by itself establish that the effect is a property of the method rather than
of that particular disorder realization. Two further independent EA disorder realizations were
therefore tested at the same $L=16$, $T=2.6$, beyond$=4$ configuration, one trained network
each -- the network-initialization axis having already been shown not to be the dominant source
of variance, a single network per realization is the appropriate allocation of computational
effort here rather than a further three-network ensemble per realization.

Table~\ref{tab:disorder} reports each realization separately, since a pooled number alone
cannot show whether one realization drives the effect. No single realization reaches
significance on its own at $n=30$: the original realization, taken as a single-network
comparison rather than the three-network pooled result above, is itself not significant
($19/30$ versus $12/30$, $p=0.120$), so this is the expected pattern at this replica count
rather than a weakening of the result. The third realization has a substantially easier
calibrated target for its particular disorder ($29/30$ and $26/30$, both near ceiling), which
is why its own test is the least informative of the three -- a ceiling effect, not evidence
against the finding.

\begin{table}[H]
\centering\small
\caption{Deep-rare EA comparison repeated across three independent disorder realizations, one
trained network each, beyond$=4$, $30$ replicas per method. Realization 12345 is the
single-network reading of the comparison in Table~\ref{tab:deep} (seed~0), not the
three-network pooled result.}
\label{tab:disorder}
\begin{tabular}{lrrr}
\toprule
Disorder realization & $\WE[I_\theta]$ & $\WE[E]$ & Fisher $p$ \\
\midrule
12345 (original) & 19/30 (63.3\%) & 12/30 (40.0\%) & 0.120 \\
23456 & 15/30 (50.0\%) & 8/30 (26.7\%) & 0.110 \\
45678 & 29/30 (96.7\%) & 26/30 (86.7\%) & 0.353 \\
\midrule
Pooled (3 realizations) & 63/90 (70.0\%) & 46/90 (51.1\%) & 0.014 \\
\bottomrule
\end{tabular}
\end{table}

\begin{figure}[H]
\centering
\includegraphics[width=\linewidth]{figs/fig_deeprare.pdf}
\caption{The learned importance coordinate reaches the deep target more reliably than the
physical energy coordinate, in the regime where direct sampling cannot compete. Error bars are
Wilson $95\%$ intervals. (a) Discovery rate at beyond$=4$ (disorder realization $12345$,
$L=16$, $T=2.6$): open diamonds are the three independently trained networks (Table~\ref{tab:deep});
$\WE[E]$ uses the balanced $90$-replica baseline ($48.9\%$; the original unbalanced $30$-replica
draw gave $40.0\%$). (b) The same comparison repeated across three independent disorder
realizations, one trained network each (Table~\ref{tab:disorder}), with the pooled result at
right.}
\label{fig:deeprare}
\end{figure}

Pooled across all three independent realizations, one network each ($90$ replicas per
method):
\begin{equation}
\WE[I_\theta]:\ 63/90\ (70.0\%) \quad\text{versus}\quad \WE[E]:\ 46/90\ (51.1\%),
\qquad p=0.014.
\label{eq:disordergen}
\end{equation}
This is a different and stronger pooling axis than the network-seed result above: rather than
asking whether the effect is robust to network initialization on one lattice, it asks whether
the effect is a property of the method across the disorder ensemble. Wall-clock cost ratios for
the two new realizations ($4.45$ and $4.97$ relative to $\WE[E]$) are consistent with every
other measurement in this study (\S\ref{sec:costdisc}), so the efficiency conclusion
generalizes together with the reliability one.

A methodological note is worth recording. A third candidate disorder realization was
attempted and abandoned before completion: its real-depth training-data collection pool
returned no states in the bulk-equilibrium reference region at all, so the self-consistency
objective's second boundary condition (\S\ref{sec:neural}) could never fire and training
collapsed toward the trivial constant solution, a failure mode the training code detects and
flags explicitly. This was verified to be a property of that particular $J_{ij}$ sample's
dynamics under the collection protocol, not a general problem: two replacement realizations
were each checked with a short collection-only diagnostic before committing to the full
run, and the one reported here was healthy by that check. Screening a new disorder realization
this way before full training is accordingly recommended practice for any extension of this
result.

\FloatBarrier

% ============================================================================
\section{Discussion}\label{sec:discussion}

The results support a direct progression from transport theory to rare-event sampling in spin
systems. The first level is mathematical: for first-passage problems, the committor is the
harmonic importance function and the Doob transform supplies the zero-variance ideal. The
birth--death demonstration verifies this limit numerically and provides a transparent reference
against which approximate importance strategies can be interpreted.

The second level is algorithmic. Weight-conserving WE resampling transfers the population-control
logic of transport weight windows without requiring the importance coordinate to be exact. Exact
enumeration on the $4\times4$ Ising model confirms that the target measure is recovered across
the tested regimes, and the $L=20$ tail calculation shows why population control becomes valuable
in practice: once the direct sample reaches its floor, the weighted population continues to
populate states with probabilities more than an order of magnitude smaller.

The third level is representational. In the EA spin glass, magnetization is not a useful
descriptor of a low-energy rare event, while energy is informative but compresses the
configuration to one scalar. A learned committor-like coordinate can use the full lattice pattern
to rank configurations by their prospects of progressing toward the target. The significant
increase from a balanced $48.9\%$ to $66.7\%$ event-positive replicas demonstrates that this
additional information can improve deep-event access. Because the effect is reproduced by three
independently trained networks on one realization, and again by pooling across three
independent disorder realizations ($51.1\%$ to $70.0\%$, $p=0.014$, \S\ref{sec:disordergen}),
the result is not tied to one isolated initialization or to one quenched-disorder sample.

\subsection{Reliability is improved; cost-normalized efficiency is not}\label{sec:costdisc}
The improvement reported in \S\ref{sec:deep} is a reliability result, and it should not be read
as a claim of computational speedup. On the spin-update cost measure, the three trained networks
give figure-of-merit ratios of $1.29$, $1.07$, and $1.04$ against $\WE[E]$
(Table~\ref{tab:deep}); a two-level bootstrap over networks and replicas gives a geometric-mean
ratio of $1.16$ with $95\%$ interval $[0.65,2.12]$, which does not exclude parity. That measure
is in any case incomplete, because it does not charge for the network evaluations performed at
every resampling step.

Direct measurement removes the need to extrapolate that omission from a proxy. All three deep-rare
networks were rerun end-to-end with wall-clock time recorded as a first-class output for every
method (not reconstructed from log timestamps), reusing the deterministic trained weights so the
underlying $\hat\pi$ estimates are unchanged. Measured wall-clock ratios of $\WE[I_\theta]$ to
$\WE[E]$ were $4.65$, $4.01$, and $4.54$ for the three networks -- consistent in direction with,
though somewhat larger than, an earlier profiling-based estimate of $2.9$--$3.4$ for a single
network obtained at a shallower configuration. The $\WE[E]$ baseline itself was also expanded
from the original fixed $30$ replicas to a balanced $90$ (\S\ref{sec:deep}); the wall-clock-fair
figure of merit uses this larger pool only to obtain a better-calibrated estimate of $\WE[E]$'s
native-replica-count precision (bootstrap-resampling $30$-replica blocks from the pool of $90$,
at $\WE[E]$'s own measured wall-clock cost for $30$ replicas), rather than charging its full
expanded wall-clock time against $\WE[I_\theta]$'s native budget, which would inflate $\WE[E]$'s
effective cost and bias the ratio. The resulting two-level-bootstrapped, wall-clock-fair
figure-of-merit ratio is
\begin{equation}
\FOM_{\rm wall}(\WE[I_\theta])\big/\FOM_{\rm wall}(\WE[E]) = 0.33,
\qquad 95\%\ \mathrm{CI}=[0.19,0.58],
\label{eq:fomwall}
\end{equation}
below parity and excluding $1$. This directly measured value agrees in direction and order of
magnitude with an earlier estimate obtained by rescaling the spin-update-cost ratio with a
profiling-based timing factor from a shallower configuration ($0.39$, $[0.22,0.72]$), which is a
useful cross-check of that method, but Eq.~\eqref{eq:fomwall} is the value that should be quoted
going forward. The corresponding value for the five-network ensemble, approximately $0.087$, is
obtained at the shallower configuration where the ensemble was characterized and is reported as
an indication of how the penalty scales with network count, not as a second deep-rare
measurement.

Table~\ref{tab:cost} consolidates the reliability and cost results so that the trade-off can be
read in one place.

\begin{table}[H]
\centering\small
\caption{Reliability and cost of the learned coordinate relative to the energy coordinate.
Wall time is relative to $\WE[E]$ under the same configuration. $\FOM_{\rm MC}$ ratios are
geometric means with two-level bootstrap $95\%$ intervals on the spin-update cost measure;
$\FOM_{\rm wall}$ ratios are directly measured (single network) or extrapolated (ensemble), see
notes. Note that the two learned-coordinate rows come from different configurations and are not
a like-for-like comparison with each other.}
\label{tab:cost}
\begin{tabular}{lccccc}
\toprule
Method & networks & rel.\ wall time & event-positive & $\FOM_{\rm MC}$ ratio &
$\FOM_{\rm wall}$ ratio \\
\midrule
$\WE[E]$\textsuperscript{a} & 0 & $1.00$ & 44/90\textsuperscript{c} & $1.00$ & $1.00$ \\
$\WE[I_\theta]$\textsuperscript{a} & 1 & $4.0$--$4.7$ & 60/90 &
$1.16$ $[0.65,2.12]$ & $0.33$ $[0.19,0.58]$ \\
$\WE[I_\theta]$\textsuperscript{b} & 5 & $10.7$--$10.9$ & --- &
$0.74$ $[0.43,1.32]$ & $0.087$ $[0.051,0.155]$ \\
\bottomrule
\end{tabular}
\\[2pt]
{\footnotesize\textsuperscript{a}Deep-rare configuration of \S\ref{sec:deep}
(beyond$=4$, 30 replicas per network), directly measured (\S\ref{sec:costdisc}).
\textsuperscript{b}Ensemble characterization at the shallower configuration
(beyond$=1$, $\tau=2$, 10 replicas per run, five independent base seeds), where every method
reaches the target in every replica so that event-positive counts are uninformative; this row
is still extrapolated from that shallower configuration, not directly measured at beyond$=4$.
\textsuperscript{c}Balanced $90$-replica baseline (\S\ref{sec:deep}); the naive fixed-30 draw
gave $12/30$.}
\end{table}

Two properties of these values should be stated plainly. The single-network deep-rare figure is
now a direct measurement rather than an extrapolation, and its interval propagates sampling
uncertainty in the $\hat\pi$ estimates on both sides of the ratio; the five-network ensemble
figure is still obtained by extrapolating a timing factor measured at a shallower configuration,
and its interval does not include uncertainty in that factor itself. The conclusion that survives
these caveats is qualitative but robust to the details: a learned coordinate can be dynamically
better and simultaneously worse per unit of computer time, and increasing neural capacity moves
the second property in the wrong direction faster than it improves the first. A dedicated
efficiency study with timing recorded as a first-class output throughout, and with inference cost
reduced by caching or batching, is the natural sequel to the present work; the single-network
figure now meets that bar, the ensemble figure does not yet.

\subsection{Relation to neighboring rare-event methods}\label{sec:relation}
The present implementation is one member of a broader splitting and path-sampling family, and
the relevant contrast is mechanistic. FFS and AMS advance trajectories through fixed interfaces
or adaptively chosen levels \citep{allen2009,valeriani2007,cerou2007}, and transition-path
sampling works directly with reactive-path ensembles, which frees it from a fixed scalar
coordinate but targets different estimators than the stationary tail probabilities studied here
\citep{chandler1978,dellago1998,bolhuis2002}. WE differs in maintaining a continuously weighted
population and resampling within bins, so that the bins influence variance but not the target
measure \citep{huber1996,zhang2010,zuckerman2017}. What the transport analogy adds is a
principled target for choosing those bins: CADIS couples an importance map to source biasing and
weight windows \citep{wagner1998,haghighat2003,wagnerfw2014}, and the present implementation
adopts the population-control half of that logic without altering the Metropolis kernel.

Against the neural rare-event literature, the distinction is where the learned object is
inserted. Flow-based and autoregressive samplers learn global equilibrium structure
\citep{wu2019,noe2019,nicoli2020,gabrie2022}; neural committor and Feynman--Kac methods target
dynamical progress toward a specific event \citep{khoo2019,li2022,strahan2023}; and active or
self-consistent committor schemes concentrate training and sampling effort in transition-relevant
regions \citep{rotskoff2022,kang2024,trizio2025}. The framework of Kim and Cai is closest in
spirit, combining a learned event-specific importance function with biased transitions and
branching \citep{kimcai2026}: in the language of \S\ref{sec:adjoint}, their importance function
plays the role of $h$, their modified transitions the role of the tilt, and their branching the
role of splitting. All of these place the learned function inside the dynamics. The present
approach places it only in the binning, which preserves unbiasedness through weight conservation
and makes coordinate quality measurable without confounding it with likelihood-ratio error.

\subsection{Scope and reproducibility}\label{sec:scope}
The principal learned-coordinate experiment is a proof of concept, and its scope should be read
literally. It uses three two-dimensional EA disorder realizations, one temperature ($T=2.6$), one
lattice size ($L=16$), one threshold depth (beyond$=4$), and a compact periodic convolutional
network. Replication in this study is across trained networks, disorder realizations, and
simulation replicas, not yet across lattice size or temperature. The result therefore establishes
that a learned configuration-space coordinate \emph{can} carry progress information beyond energy
on a frustrated landscape, and does so reproducibly across network initializations
(\S\ref{sec:deep}) and across three independent disorder realizations (\S\ref{sec:disordergen});
it does not establish generality across lattice size or temperature, and the three-realization
sample is not large enough to characterize how the effect's magnitude varies with disorder --
one of the three realizations tested presented a substantially easier calibrated target than the
other two. A larger disorder sample, to tighten the pooled confidence interval and better
characterize this variation, and extensions in $L$ and $T$, are the natural next steps, and are
more expensive than the disorder-axis test already completed. Exact enumeration, the solvable
birth--death reference, fixed-seed reproducibility checks, and explicit reporting of zero-hit
replicas provide complementary safeguards for the numerical interpretation. One further caution
follows from the project's own experience: because outcomes vary appreciably between training
initializations at this scale, a single training run is not evidence about a learned coordinate
in either direction, and seed-resolved reporting should be treated as a requirement rather than a
refinement. All learned-coordinate results reported here are accordingly stated across
independently trained networks. A second, disorder-specific caution also follows from this
project's experience: not every disorder realization is trainable under the self-consistency
objective's collection protocol, since the second boundary condition can fail to fire for a
particular $J_{ij}$ sample's dynamics (\S\ref{sec:disordergen}); a brief collection-only
diagnostic before committing to full training is accordingly recommended when extending this
result to further realizations.

The next natural extension is to test the same learned-coordinate construction across lattice
sizes and lower-temperature regimes where the separation between common and rare configurations
becomes stronger, and across a larger sample of disorder realizations. Architectures that encode
spin symmetries or local frustration explicitly may further improve the quality of the learned
importance map. The same framework can also be coupled to adaptive bins or milestone definitions,
provided the underlying split/merge step remains weight conserving.

\FloatBarrier

% ============================================================================
\section{Conclusions}\label{sec:conclusions}
Adjoint transport theory provides a useful organizing principle for rare-event sampling in lattice
spin systems. The committor supplies the event-specific importance function, the exact Doob
transform defines the zero-variance ideal, and weight-conserving population control provides a
practical way to exploit approximate importance without changing the target measure. These
elements were validated sequentially from a solvable birth--death process to exact enumeration of
the Ising model.

At system scale, WE reaches an $L=20$ Ising magnetization-tail bin approximately $24$ times
below the direct-sampling floor near $5.9\times10^{-6}$ at matched spin-update cost, returning a
(noisy) estimate of $2.4\times10^{-7}$ where direct sampling returns none. In the frustrated
EA model, the learned configuration-space importance coordinate produces the main new result:
pooled across three independent disorder realizations with one trained network each, it yields
63 event-positive replicas out of 90 against 46/90 for the physical energy coordinate
($70.0\%$ versus $51.1\%$, Fisher exact $p=0.014$); a complementary pooling of three
independently trained networks on a single realization against a matched-size baseline gives
$66.7\%$ versus $48.9\%$ ($p=0.023$). The result shows that a learned progress coordinate can
extract dynamical information that is not fully represented by a scalar physical observable, and
that this advantage is reproducible both across network initializations and across independent
disorder realizations -- a property of the method rather than of one particular quenched-disorder
sample. A larger disorder sample would further tighten this estimate and is the natural next
step.

That advantage is presently one of reliability rather than efficiency. Once repeated network
inference is directly measured and charged against the result, the learned coordinate sits below
the energy coordinate in wall-clock figure of merit by a factor of approximately three for a
single network ($0.33$, $95\%$ CI $[0.19,0.58]$, excluding parity) and by an order of magnitude
for a five-network ensemble. Reducing that inference cost, rather than improving the coordinate
further, is therefore the binding constraint on making learned importance practical at this
scale, and is the specific target of the follow-up study.

The resulting framework connects classical transport variance reduction, WE population sampling,
and neural importance learning in a single validated workflow. It is directly extensible to larger
spin systems and to other rare-event problems in which a physically obvious one-dimensional
coordinate is insufficient.

% ============================================================================
\section*{Data and Code Availability}
All code, input data, and result files required to reproduce the results reported here are
publicly archived \citep{dataset}. The deposit contains the Weighted Ensemble implementation and
production drivers, the training code for both objectives, the quenched disorder realizations used
for every Edwards--Anderson result, per-replica $\hat\pi$ arrays for every table and figure, the
exact-enumeration and birth--death validation outputs, the trained network weights for the three
within-disorder deep-rare networks and the two additional disorder-realization networks of
\S\ref{sec:disordergen}, the raw wall-clock logs underlying the cost analysis of
\S\ref{sec:costdisc},
and a manifest mapping each published number to the file that produced it. An environment
specification pins package versions, and all published results are CPU-deterministic under the
recorded seeds. Results not reported in the main text, including the seed sweep and ensemble runs
of \S\ref{sec:costdisc}, are retained in full in the same deposit and described in the
Supplementary Material.

\section*{Acknowledgements}
The author gratefully acknowledges the support of Fahad Bin Sultan University, Tabuk, Saudi Arabia,
including access to the computational environment used for the numerical experiments.

\printcredits

% ============================================================================
\begin{thebibliography}{99}

\bibitem{ising1925} E. Ising, Beitrag zur Theorie des Ferromagnetismus, Z. Phys. 31 (1925)
253--258.
\bibitem{onsager1944} L. Onsager, Crystal statistics. I. A two-dimensional model with an
order-disorder transition, Phys. Rev. 65 (1944) 117--149.
\bibitem{edwards1975} S. F. Edwards, P. W. Anderson, Theory of spin glasses, J. Phys. F: Met.
Phys. 5 (1975) 965--974.
\bibitem{binder1986} K. Binder, A. P. Young, Spin glasses: experimental facts, theoretical
concepts, and open questions, Rev. Mod. Phys. 58 (1986) 801--976.
\bibitem{metropolis1953} N. Metropolis, A. W. Rosenbluth, M. N. Rosenbluth, A. H. Teller,
E. Teller, Equation of state calculations by fast computing machines, J. Chem. Phys. 21 (1953)
1087--1092.
\bibitem{kahn1951} H. Kahn, T. E. Harris, Estimation of particle transmission by random sampling,
Natl. Bur. Stand. Appl. Math. Ser. 12 (1951) 27--30.
\bibitem{kalos2008} M. H. Kalos, P. A. Whitlock, Monte Carlo Methods, 2nd ed., Wiley-VCH,
Weinheim, 2008.
\bibitem{lewins1965} J. Lewins, Importance: The Adjoint Function, Pergamon Press, Oxford, 1965.
\bibitem{spanier1969} J. Spanier, E. M. Gelbard, Monte Carlo Principles and Neutron Transport
Problems, Addison-Wesley, Reading, MA, 1969.
\bibitem{bell1970} G. I. Bell, S. Glasstone, Nuclear Reactor Theory, Van Nostrand Reinhold,
New York, 1970.
\bibitem{wagner1998} J. C. Wagner, A. Haghighat, Automated variance reduction of Monte Carlo
shielding calculations using the discrete ordinates adjoint function, Nucl. Sci. Eng. 128 (1998)
186--208.
\bibitem{haghighat2003} A. Haghighat, J. C. Wagner, Monte Carlo variance reduction with
deterministic importance functions, Prog. Nucl. Energy 42 (2003) 25--53.
\bibitem{wagnerfw2014} J. C. Wagner, D. E. Peplow, S. W. Mosher, FW-CADIS method for global and
regional variance reduction of Monte Carlo radiation transport calculations, Nucl. Sci. Eng. 176
(2014) 37--57.
\bibitem{huber1996} G. A. Huber, S. Kim, Weighted-ensemble Brownian dynamics simulations for
protein association reactions, Biophys. J. 70 (1996) 97--110.
\bibitem{zhang2010} B. W. Zhang, D. M. Zuckerman, D. Jasnow, The weighted ensemble path sampling
method is statistically exact for a broad class of stochastic processes and binning procedures,
J. Chem. Phys. 132 (2010) 054107.
\bibitem{bhatt2010} D. Bhatt, B. W. Zhang, D. M. Zuckerman, Steady-state simulations using
weighted ensemble path sampling, J. Chem. Phys. 133 (2010) 014110.
\bibitem{zuckerman2017} D. M. Zuckerman, L. T. Chong, Weighted ensemble simulation: review of
methodology, applications, and software, Annu. Rev. Biophys. 46 (2017) 43--57.
\bibitem{russo2022} J. D. Russo, et al., WESTPA 2.0: high-performance upgrades for weighted
ensemble simulations and analysis of longer-timescale applications, J. Chem. Theory Comput. 18
(2022) 638--649.
\bibitem{allen2009} R. J. Allen, C. Valeriani, P. R. ten Wolde, Forward flux sampling for rare
event simulations, J. Phys.: Condens. Matter 21 (2009) 463102.
\bibitem{valeriani2007} C. Valeriani, R. J. Allen, M. J. Morelli, D. Frenkel, P. R. ten Wolde,
Computing stationary distributions in equilibrium and nonequilibrium systems with forward flux
sampling, J. Chem. Phys. 127 (2007) 114109.
\bibitem{cerou2007} F. C\'erou, A. Guyader, Adaptive multilevel splitting for rare event analysis,
Stoch. Anal. Appl. 25 (2007) 417--443.
\bibitem{dellago1998} C. Dellago, P. G. Bolhuis, D. Chandler, Efficient transition path sampling:
application to Lennard-Jones cluster rearrangements, J. Chem. Phys. 108 (1998) 9236--9245.
\bibitem{bolhuis2002} P. G. Bolhuis, D. Chandler, C. Dellago, P. L. Geissler, Transition path
sampling: throwing ropes over rough mountain passes, in the dark, Annu. Rev. Phys. Chem. 53 (2002)
291--318.
\bibitem{e2006} W. E, E. Vanden-Eijnden, Towards a theory of transition paths, J. Stat. Phys. 123
(2006) 503--523.
\bibitem{metzner2009} P. Metzner, C. Sch\"utte, E. Vanden-Eijnden, Transition path theory for
Markov jump processes, Multiscale Model. Simul. 7 (2009) 1192--1219.
\bibitem{vandeneijnden2010} E. Vanden-Eijnden, Transition path theory, in: M. Ferrario,
G. Ciccotti, K. Binder (Eds.), Computer Simulations in Condensed Matter Systems: From Materials to
Chemical Biology, Vol. 1, Springer, 2010, pp. 453--493.
\bibitem{doob1984} J. L. Doob, Classical Potential Theory and Its Probabilistic Counterpart,
Springer, New York, 1984.
\bibitem{touchette2009} H. Touchette, The large deviation approach to statistical mechanics,
Phys. Rep. 478 (2009) 1--69.
\bibitem{chetrite2015} R. Chetrite, H. Touchette, Nonequilibrium Markov processes conditioned on
large deviations, Ann. Henri Poincar\'e 16 (2015) 2005--2057.
\bibitem{chandler1978} D. Chandler, Statistical mechanics of isomerization dynamics in liquids and
the transition state approximation, J. Chem. Phys. 68 (1978) 2959--2970.
\bibitem{wu2019} D. Wu, L. Wang, P. Zhang, Solving statistical mechanics using variational
autoregressive networks, Phys. Rev. Lett. 122 (2019) 080602.
\bibitem{noe2019} F. No\'e, S. Olsson, J. K\"ohler, H. Wu, Boltzmann generators: sampling
equilibrium states of many-body systems with deep learning, Science 365 (2019) eaaw1147.
\bibitem{nicoli2020} K. A. Nicoli, et al., Asymptotically unbiased estimation of physical
observables with neural samplers, Phys. Rev. E 101 (2020) 023304.
\bibitem{gabrie2022} M. Gabri\'e, G. M. Rotskoff, E. Vanden-Eijnden, Adaptive Monte Carlo augmented
with normalizing flows, Proc. Natl. Acad. Sci. U.S.A. 119 (2022) e2109420119.
\bibitem{carleo2017} G. Carleo, M. Troyer, Solving the quantum many-body problem with artificial
neural networks, Science 355 (2017) 602--606.
\bibitem{carrasquilla2017} J. Carrasquilla, R. G. Melko, Machine learning phases of matter,
Nat. Phys. 13 (2017) 431--434.
\bibitem{khoo2019} Y. Khoo, J. Lu, L. Ying, Solving for high-dimensional committor functions using
artificial neural networks, Res. Math. Sci. 6 (2019) 1.
\bibitem{li2022} H. Li, Y. Khoo, Y. Ren, L. Ying, A semigroup method for high dimensional committor
functions based on neural network, Proc. Mach. Learn. Res. 145 (2022) 598--618.
\bibitem{rotskoff2022} G. M. Rotskoff, A. R. Mitchell, E. Vanden-Eijnden, Active importance
sampling for variational objectives dominated by rare events: consequences for optimization and
generalization, Proc. Mach. Learn. Res. 145 (2022) 757--780.
\bibitem{strahan2023} J. Strahan, J. Finkel, A. R. Dinner, J. Weare, Predicting rare events using
neural networks and short-trajectory data, J. Comput. Phys. 488 (2023) 112152.
\bibitem{kang2024} P. Kang, E. Trizio, M. Parrinello, Computing the committor with the committor to
study the transition state ensemble, Nat. Comput. Sci. 4 (2024) 451--460.
\bibitem{trizio2025} E. Trizio, P. Kang, M. Parrinello, Everything everywhere all at once: a
probability-based enhanced sampling approach to rare events, Nat. Comput. Sci. (2025).
\bibitem{hua2024} X. Hua, R. Ahmad, J. Blanchet, W. Cai, Accelerated sampling of rare events using
a neural network bias potential, arXiv:2401.06936 (2024).
\bibitem{kimcai2026} M. Kim, W. Cai, Accelerated Markov Chain Monte Carlo simulation via neural
network-driven importance sampling, arXiv:2602.12294 (2026).
\bibitem{zhao2025} Y. Zhao, et al., Nonequilibrium statistical mechanics revealed by Doob
$h$-transforms and variational neural representations, Phys. Rev. E 111 (2025) 034120.
\bibitem{fisher1922} R. A. Fisher, On the interpretation of $\chi^2$ from contingency tables,
and the calculation of $P$, J. R. Stat. Soc. 85 (1922) 87--94.
\bibitem{efron1979} B. Efron, Bootstrap methods: another look at the jackknife, Ann. Statist.
7 (1979) 1--26.
\bibitem{dataset} M. Y. Alzaatreh, Code and data for ``Transport-inspired rare-event sampling in
spin systems'', Zenodo (2026). doi:10.5281/zenodo.XXXXXXX.

\end{thebibliography}

\end{document}
